\chapter{Zeros and Primes in Short Intervals}
\begin{center}\itshape
Is there always a prime between $x$ and $x+\sqrt x$? Nobody knows, and the Riemann hypothesis would only just fail to tell.\\
But one does not need every zero on the critical line: it is enough that few of them stray far from it.\\
Counting the strays is the art of this chapter, and its prize is a prime in every interval $\left]x,x+x^{0.76}\right]$.
\end{center}

Between $436\,273\,009$ and the next prime, $436\,273\,291$, there are $281$ composite numbers in a row. Below a billion no gap is longer, and the record gaps grow so slowly that $(\log p)^2$, about $396$ there, sits above every record gap that starts at a prime $p\ge23$. Harald Cram\'er
conjectured in 1936 that this is the truth: gaps of size at most a constant times $(\log p)^2$. What has been proved is far weaker. The Riemann hypothesis would give a prime in every interval $\left]x,x+\sqrt x\log^3x\right]$; this is Cram\'er's theorem of 1920, an exercise in Volume~5. Without
any hypothesis, the prime number theorem in its classical form only finds primes in intervals $\left]x,x+\varepsilon x\right]$: for every fixed $\varepsilon>0$ and $x$ large, but the length is still proportional to $x$.

In 1930 Guido Hoheisel (1894--1968) broke the barrier: there is a constant $\theta<1$ such that every interval $\left]x,x+x^\theta\right]$ contains about $x^\theta/\log x$ primes, for $x$ large. His $\theta$ was $1-\frac1{33000}$, and the method mattered more
than the number. Hoheisel did not need the zeros of $\zeta$ to lie on the critical line; he needed that \emph{few} zeros lie far to its right. That idea goes back to 1914, when Harald Bohr and Edmund Landau proved that almost all the zeros, in proportion, lie within any distance $\delta$
of the line, and to Fritz Carlson, who in 1920 counted the exceptions: at most about $T^{4\sigma(1-\sigma)}$ zeros up to height $T$ have real part $\ge\sigma$. Albert Ingham turned Hoheisel's argument into a clean machine in 1937: a bound $T^{A(1-\sigma)}$ for the number of zeros to the right of $\sigma$ gives primes in all intervals of length $x^\theta$ with $\theta>1-\frac1A$.
Carlson's bound gives $A=4$, so $\theta>\frac34$; Ingham's own gives $\frac58$, and the race has continued ever since.

This chapter proves all of this, except one zero-free region which is quoted. The tools are humble: a mean value theorem for Dirichlet polynomials, Jensen's formula from Volume~5, and a \emph{mollifier}, a short Dirichlet polynomial that cancels $\zeta$ on average.

\begin{reminder}[from Volume 5]
$\psi(x)=\sum_{n\le x}\Lambda(n)$, and $\psi_0(x)$ is the same with the last term halved when $x$ is a prime power. The zeros $\rho=\beta+i\gamma$ of $\zeta$ in the critical strip are symmetric under $\rho\mapsto\overline\rho$ and $\rho\mapsto1-\overline\rho$. $N(T)$, the number of zeros with $0<\gamma\le T$, is
$\frac T{2\pi}\log\frac T{2\pi e}+O(\log T)$ (Theorem~6.9), and there are $O(\log T)$ zeros with $T<\gamma\le T+1$ (Corollary~6.10). No zero has $\beta=1$ (Theorem~6.12). The explicit formula (Theorem~7.5): for $x\ge2$, $T\ge3$,
\[
\begin{aligned}\psi_0(x)&=x-\sum_{|\gamma|<T}\frac{x^\rho}\rho-\log2\pi-\frac12\log\Bigl(1-\frac1{x^2}\Bigr)\eqbr+O\Bigl(\frac{x\log^2(xT)}T+\log x\Bigr).\end{aligned}
\]
Jensen's formula (Theorem~6.3): if $h$ is holomorphic on a neighbourhood of the closed disc $|z-a|\le r$, $h(a)\ne0$ and $h$ has no zero on $|z-a|=r$, then $\log|h(a)|+\sum\log\frac r{|z_j-a|}=\frac1{2\pi}\int_0^{2\pi}\log|h(a+re^{i\theta})|\dd\theta$, the sum over the zeros $z_j$ in $|z-a|<r$ with multiplicity.
And for $\operatorname{Re}s>0$, $s\ne1$, $N\ge1$ an integer (Volume~5, Exercise~6.6):
\[
\zeta(s)=\sum_{n\le N}n^{-s}+\frac{N^{1-s}}{s-1}-s\int_N^\infty\{u\}u^{-s-1}\dd u.
\]
\end{reminder}

\section{From zeros to primes in short intervals}

Throughout, $\rho=\beta+i\gamma$ runs over the non-trivial zeros of $\zeta$, with multiplicity, and for $\frac12\le\sigma\le1$ and $T>0$,
\[
N(\sigma,T)=\#\{\rho:\ \beta\ge\sigma,\ 0<\gamma\le T\}.
\]
Why should this count control primes in short intervals? The explicit formula writes $\psi(x+h)-\psi(x)$ as $h$ minus a sum over zeros, and each zero contributes
\[
\frac{(x+h)^\rho-x^\rho}\rho=\int_x^{x+h}u^{\rho-1}\dd u,\qquad\text{of modulus at most }hx^{\beta-1}.
\]
So the zeros enter through $x^{\beta-1}$: a zero on the critical line costs $hx^{-1/2}$, a zero near $\operatorname{Re}s=1$ costs almost $h$. Only the zeros far to the right are dangerous, and the question is whether there are few enough of them to be paid for. Near $\operatorname{Re}s=1$ itself we need a zero-free region, and
the classical one of Volume~5 is too thin. We quote a better one.

\begin{theorem}[Vinogradov, Korobov, 1958; quoted]\label{thm:vk}
There is an absolute constant $c_0>0$ such that $\zeta(\sigma+it)\ne0$ whenever $|t|\ge3$ and
\[
\sigma\ge1-\frac{c_0}{(\log|t|)^{2/3}(\log\log|t|)^{1/3}}.
\]
\end{theorem}

Its proof rests on Vinogradov's method for exponential sums, which Chapter~5 introduces; Kevin Ford made it explicit in 2002, with $c_0=\frac1{57.54}$. We use it only in the next theorem, and only through one property: the width $\eta(T)$ of the region at height $T$ satisfies $\eta(T)\log T/\log\log T\to\infty$.

\begin{theorem}[Ingham, 1937]\label{thm:ingham}
Suppose that there are constants $A\ge2$, $B\ge0$ and $C$ with
\[
N(\sigma,T)\le CT^{A(1-\sigma)}(\log T)^B\qquad\text{for all }\tfrac12\le\sigma\le1,\ T\ge2.
\]
Then for every fixed $\theta$ with $1-\frac1A<\theta\le1$, $\psi(x+x^\theta)-\psi(x)\sim x^\theta$ as $x\to\infty$.
\end{theorem}

\begin{proof}
Let $h=x^\theta$ and $T=x^{1-\theta}(\log x)^3$; for $x$ large, $3\le T\le x$ and $\log(xT)\le2\log x$.

\emph{The explicit formula.} Apply it at $x$ and at $x+h\le2x$, with the same $T$, and subtract. The terms $\log2\pi$ cancel, the terms $\frac12\log(1-y^{-2})$ are $O(x^{-2})$, and $|\psi_0(y)-\psi(y)|\le\frac12\log y$. So
\[
\psi(x+h)-\psi(x)=h-\sum_{|\gamma|<T}\frac{(x+h)^\rho-x^\rho}\rho+O\Bigl(\frac{x\log^2x}T+\log x\Bigr).
\]
The error is $O(\frac{x^\theta}{\log x})=o(h)$, and by the remark above the sum over zeros has modulus at most $hS$, where $S=\sum_{|\gamma|<T}x^{\beta-1}$. It remains to show that $S\to0$.

\emph{Low zeros.} There are finitely many zeros with $|\gamma|<3$, and each has $\beta<1$; if $\eta_0>0$ is the least of their $1-\beta$, they contribute $O(x^{-\eta_0})\to0$. The zeros with $3\le|\gamma|<T$ come in conjugate pairs, so they contribute twice $S'=\sum_{3\le\gamma<T}x^{\beta-1}$.

\emph{Layer cake.} For a zero with $\beta\le\frac12$, $x^{\beta-1}\le x^{-1/2}$; for $\beta>\frac12$, $x^{\beta-1}=x^{-1/2}+\log x\int_{1/2}^\beta x^{\sigma-1}\dd\sigma$. Summing, and exchanging sum and integral (finitely many terms),
\[
S'\le N(T)x^{-1/2}+\log x\int_{1/2}^1x^{\sigma-1}N(\sigma,T)\dd\sigma.
\]
The first term is $O(Tx^{-1/2}\log T)$, which tends to $0$ because $1-\theta<\frac1A\le\frac12$.

\emph{The zero-free region.} Let $\eta=\frac{c_0}{(\log T)^{2/3}(\log\log T)^{1/3}}$. The function $t\mapsto\frac{c_0}{(\log t)^{2/3}(\log\log t)^{1/3}}$ decreases on $[3,\infty[$, so Theorem~\ref{thm:vk} shows that no zero with $3\le\gamma<T$ has $\beta\ge1-\eta$. Hence the integral may stop at $1-\eta$, and with the hypothesis,
\[
\log x\int_{1/2}^{1-\eta}x^{\sigma-1}N(\sigma,T)\dd\sigma\le C(\log T)^B\log x\int_{1/2}^{1-\eta}\Bigl(\frac{T^A}x\Bigr)^{1-\sigma}\dd\sigma.
\]
Now $\frac{T^A}x=x^{-\delta}(\log x)^{3A}$ with $\delta=1-A(1-\theta)>0$, so $\frac{T^A}x\le y:=x^{-\delta/2}$ for $x$ large. Substituting $u=1-\sigma$, $\int_\eta^{1/2}y^u\dd u\le\frac{y^\eta}{\log(1/y)}=\frac{2x^{-\delta\eta/2}}{\delta\log x}$. Since $\log T\le\log x$, $\eta\log x\ge\frac{c_0(\log x)^{1/3}}{(\log\log x)^{1/3}}$, and the right side above is at most
\[
\frac{2C}\delta(\log x)^B\exp\Bigl(-\frac{\delta c_0}2\cdot\frac{(\log x)^{1/3}}{(\log\log x)^{1/3}}\Bigr)\longrightarrow0,
\]
because $(\log x)^{1/3}(\log\log x)^{-1/3}$ grows faster than any multiple of $\log\log x$. So $S\to0$, and $\psi(x+h)-\psi(x)=h+o(h)$.
\end{proof}

\begin{think}
Where in this proof would the classical zero-free region $\sigma\ge1-\frac c{\log t}$ of Volume~5 fail? Try it: what does $x^{-\delta\eta/2}$ become when $\eta=\frac c{\log T}$ and $T$ is a power of $x$?
\end{think}

With $\eta=\frac c{\log T}$ and $\log T\asymp\log x$, the factor $x^{-\delta\eta/2}$ is a constant, and nothing pays for $(\log x)^B$. The density bound and the zero-free region must work together: one controls the zeros at distance more than $\eta$ from $\operatorname{Re}s=1$, the other forbids the closer ones, and the logarithms in between need the
region to be a little wider than $\frac1{\log T}$. So the rest of the chapter has a single goal: a density bound with some $A$. Before counting zeros we need a way to measure the size of Dirichlet polynomials on average.

\section{Mean values of Dirichlet polynomials}

A \emph{Dirichlet polynomial} is a finite sum $F(t)=\sum_na_nn^{-it}$. Expanding $|F(t)|^2=\sum_{m,n}a_m\overline{a_n}(n/m)^{it}$, the diagonal $m=n$ gives $\sum|a_n|^2$ for every $t$, and the other terms oscillate. Over a long interval they should average out.

\begin{lemma}[a mean value theorem]\label{lem:mvt}
Let $M\ge1$, let $a_n$ be complex numbers for $M<n\le2M$, $T_0\in\R$ and $T>0$. Then
\[
\int_{T_0}^{T_0+T}\Bigl|\sum_{M<n\le2M}a_nn^{-it}\Bigr|^2\dd t\le\bigl(T+8M(1+\log M)\bigr)\sum_{M<n\le2M}|a_n|^2.
\]
\end{lemma}
\begin{proof}
Expand the square and integrate term by term. The diagonal gives $T\sum|a_n|^2$. For $m\ne n$, $\bigl|\int_{T_0}^{T_0+T}e^{it\log(n/m)}\dd t\bigr|\le\frac2{|\log(n/m)|}$, and by $|a_m\overline{a_n}|\le\frac12(|a_m|^2+|a_n|^2)$ and symmetry the off-diagonal part is at most
\[
\sum_{m\ne n}\frac{|a_m|^2+|a_n|^2}{|\log(n/m)|}=2\sum_m|a_m|^2\sum_{n\ne m}\frac1{|\log(n/m)|}.
\]
For $0<a<b$, $\log\frac ba\ge\frac{b-a}b$ (the mean value theorem), so $|\log\frac nm|\ge\frac{|n-m|}{2M}$ for $m,n$ in $\left]M,2M\right]$. For fixed $m$, each value $1\le d<M$ of $|n-m|$ is taken at most twice, and $\sum_{n\ne m}\frac1{|\log(n/m)|}\le2M\cdot2\sum_{d<M}\frac1d\le4M(1+\log M)$.
\end{proof}

The bound is sharp up to the logarithm: the diagonal term $T\sum|a_n|^2$ is always there, and Exercise~\ref{ex:mvtsharp} shows that the term $M\sum|a_n|^2$ is needed too. Hugh Montgomery and Robert Vaughan proved in 1974 the exact form $\sum(T+O(n))|a_n|^2$, with Hilbert's inequality; the logarithm costs us nothing here.
The coefficients we meet are bounded by the divisor function $d(n)$, and we need their mean square.

\begin{lemma}\label{lem:dsq}
For $y\ge1$, $\sum_{n\le y}\frac{d(n)^2}n\le(1+\log y)^4$; hence $\sum_{M<n\le2M}d(n)^2\le2M(1+\log2M)^4$ for $M\ge1$.
\end{lemma}
\begin{proof}
Let $d_4(n)$ be the number of ways to write $n=abce$ with positive integers. Both $d^2$ and $d_4$ are multiplicative, with $d(p^j)^2=(j+1)^2$ and $d_4(p^j)=\binom{j+3}3$ (place $j$ factors $p$ in four boxes), and $(j+1)^2\le\frac{(j+1)(j+2)(j+3)}6$ because $6(j+1)\le(j+2)(j+3)$. So $d(n)^2\le d_4(n)$, and
\[
\sum_{n\le y}\frac{d_4(n)}n=\sum_{\substack{a,b,c,e\\abce\le y}}\frac1{abce}\le\Bigl(\sum_{a\le y}\frac1a\Bigr)^4\le(1+\log y)^4.
\]
For the second claim, $\sum_{M<n\le2M}d(n)^2\le2M\sum_{n\le2M}\frac{d(n)^2}n$.
\end{proof}

\section{The zeta function as a Dirichlet polynomial}

To use mean values we must replace $\zeta(s)$ by a Dirichlet polynomial. The formula of the reminder does it with an error of size $|s|N^{-\sigma}$, which forces $N$ to be much larger than $|t|$. Integrating by parts again and again lowers the error, with periodic functions in the role of Bernoulli polynomials.

Let $P_1(u)=\{u\}-\frac12$. If $P_j$ is $1$-periodic, integrable, with $\int_0^1P_j=0$ and $|P_j|\le\frac12$, let $P_{j+1}$ be the antiderivative $u\mapsto\int_0^uP_j$ plus the constant that makes its mean over $[0,1]$ zero. It is continuous, and $1$-periodic because $\int_u^{u+1}P_j=0$. Having mean zero, it vanishes at some $u_0$, and for $u_0\le u\le u_0+1$,
$|P_{j+1}(u)|=|\int_{u_0}^uP_j|\le\int_{u_0}^{u_0+1}|P_j|\le\frac12$; by periodicity $|P_{j+1}|\le\frac12$ everywhere. So all the $P_j$ are $1$-periodic, of mean zero, and bounded by $\frac12$. (In fact $P_2=\frac{\{u\}^2-\{u\}}2+\frac1{12}$; see Exercise~\ref{ex:pk}.)

\begin{lemma}[$\zeta$ as a short sum]\label{lem:em}
Let $k\ge2$ and $N\ge1$ be integers, and $s=\sigma+it$ with $\sigma>0$, $s\ne1$. If $N\ge(|s|+k)^{k/(k-1)}$, then
\[
\zeta(s)=\sum_{n\le N}n^{-s}+\frac{N^{1-s}}{s-1}+E,\qquad|E|\le\frac{k+1}2N^{-\sigma}.
\]
\end{lemma}
\begin{proof}
Write $\{u\}=\frac12+P_1(u)$ in the formula of the reminder. Since $s\int_N^\infty\frac12u^{-s-1}\dd u$ equals $\frac12N^{-s}$, we get
\[
E=-\frac12N^{-s}-s\int_N^\infty P_1(u)u^{-s-1}\dd u.
\]
For $j\ge1$, integration by parts with $P_{j+1}'=P_j$ (except at integers when $j=1$, where $P_2$ is still continuous) gives
\[
\begin{aligned}\int_N^\infty P_j(u)u^{-s-j}\dd u&=-P_{j+1}(N)N^{-s-j}\eqbr+(s+j)\int_N^\infty P_{j+1}(u)u^{-s-j-1}\dd u,\end{aligned}
\]
the boundary term at infinity vanishing since $\sigma>0$. Starting from $j=1$ and repeating $k-1$ times,
\[
\begin{aligned}s\int_N^\infty P_1u^{-s-1}\dd u&=-\sum_{j=1}^{k-1}s(s+1)\cdots(s+j-1)P_{j+1}(N)N^{-s-j}\\&\quad+s(s+1)\cdots(s+k-1)\int_N^\infty P_k(u)u^{-s-k}\dd u.\end{aligned}
\]
Each product of $j\le k$ factors has modulus at most $(|s|+k)^j$, and $|P_j|\le\frac12$. The hypothesis says $(|s|+k)^k\le N^{k-1}$, so also $|s|+k\le N$. Hence each of the $k-1$ boundary terms is at most $\frac12N^{-\sigma}$, and the last term at most
$\frac12(|s|+k)^k\int_N^\infty u^{-\sigma-k}\dd u\le\frac12(|s|+k)^kN^{1-k-\sigma}\le\frac12N^{-\sigma}$. With the term $\frac12N^{-s}$, $|E|\le\frac{k+1}2N^{-\sigma}$.
\end{proof}

So for every $\varepsilon>0$, $\zeta(s)$ is a Dirichlet polynomial of length $|t|^{1+\varepsilon}$, up to an error of the size of its last term. (Chapter~14 brings the length down to $|t|$, and the \emph{approximate functional equation} quoted there needs only two sums of length $\sqrt{|t|/2\pi}$.)

\section{Mollifiers and Carlson's density theorem}

How can one count zeros? Jensen's formula counts the zeros of a holomorphic function $h$ in a disc by the average of $\log|h|$ on the boundary. If $h$ is close to $1$ on average, few zeros. The function $\zeta$ itself is not close to $1$; the trick, due to Bohr and Landau, is to multiply it by an approximation of $\frac1\zeta$,
\[
M_X(s)=\sum_{d\le X}\frac{\mu(d)}{d^s},
\]
the truncated Dirichlet series of $\frac1\zeta$, and to look at $f=\zeta M_X-1$. In $\operatorname{Re}s>1$,
\[
\zeta(s)M_X(s)=\sum_{n\ge1}\frac{b_n}{n^s},\qquad b_n=\sum_{d\mid n,\ d\le X}\mu(d),
\]
and $b_1=1$, $b_n=\sum_{d\mid n}\mu(d)=0$ for $1<n\le X$, $|b_n|\le d(n)$: the product starts $1+(\text{terms with }n>X)$. A function that makes $\zeta$ closer to $1$ in this way is called a \emph{mollifier}. Every zero of $\zeta$ is a zero of $\zeta M_X$, hence of
\[
h=1-f^2=\zeta M_X(2-\zeta M_X),
\]
with at least the same multiplicity, and $\log|h|\le\log(1+|f|^2)\le|f|^2$. So it suffices to show that $f$ is small in mean square.

\begin{proposition}[the mollified mean square]\label{prop:meansq}
Let $0<\varepsilon\le1$. There are $C(\varepsilon)$ and $T_1(\varepsilon)$ such that for $T\ge T_1(\varepsilon)$, $1\le X\le T$, $N=\lceil(5T)^{1+\varepsilon/2}\rceil$, $f=\zeta M_X-1$ and every $\sigma'\in[\frac12,T]$,
\[
\int_{T/2}^{3T}|f(\sigma'+it)|^2\dd t\le C(\varepsilon)(\log T)^7\bigl(TX^{1-2\sigma'}+(NX)^{2-2\sigma'}+1\bigr).
\]
\end{proposition}
\begin{proof}
Let $k=\lceil\frac2\varepsilon\rceil+1$, so $\frac k{k-1}\le1+\frac\varepsilon2$, and take $T\ge k$. For $s=\sigma'+it$ with $\frac T2\le t\le3T$, $|s|+k\le4T+k\le5T$, so $N\ge(|s|+k)^{k/(k-1)}$ and Lemma~\ref{lem:em} applies. Multiplying by $M_X(s)$,
\[
f(s)=\sum_{X<m\le NX}b'_mm^{-s}+R(s),\qquad b'_m=\sum_{\substack{d\mid m,\ d\le X\\m/d\le N}}\mu(d),\qquad R=\Bigl(\frac{N^{1-s}}{s-1}+E\Bigr)M_X(s).
\]
Indeed $\sum_{n\le N}n^{-s}M_X(s)=\sum_{m\le NX}b'_mm^{-s}$, and for $m\le X$ (so $m\le N$) the condition $m/d\le N$ is automatic: $b'_m=b_m$, which is $1$ for $m=1$ and $0$ for $1<m\le X$. Also $|b'_m|\le d(m)$.

\emph{The remainder.} $|s-1|\ge t\ge\frac T2$ and $N\ge T$, so $|\frac{N^{1-s}}{s-1}|\le\frac{2N^{1-\sigma'}}T$ and $|E|\le\frac{k+1}2N^{-\sigma'}\le\frac{k+1}{2T}N^{1-\sigma'}$. And $|M_X(s)|\le\sum_{d\le X}d^{-1/2}\le2\sqrt X\le2\sqrt T$. So $|R|\le(k+5)N^{1-\sigma'}T^{-1/2}$ and $\int_{T/2}^{3T}|R|^2\le\frac52(k+5)^2N^{2-2\sigma'}$.

\emph{The polynomial.} Cut $\left]X,NX\right]$ into $J=\lceil\log_2N\rceil$ blocks $\left]M_j,2M_j\right]$, $M_j=2^jX$, and call $P_j$ the corresponding parts of the sum. By Cauchy--Schwarz, $|\sum_jP_j|^2\le J\sum_j|P_j|^2$. Lemma~\ref{lem:mvt} on an interval of length $\frac52T$, with $a_m=b'_mm^{-\sigma'}$, and Lemma~\ref{lem:dsq} give, with $L=1+\log(2NX)$,
\[
\int_{T/2}^{3T}|P_j|^2\le\bigl(3T+8M_jL\bigr)M_j^{-2\sigma'}\cdot2M_jL^4\le2L^5\bigl(3TM_j^{1-2\sigma'}+8M_j^{2-2\sigma'}\bigr).
\]
Since $\sigma'\ge\frac12$ and $X\le M_j\le NX$, $M_j^{1-2\sigma'}\le X^{1-2\sigma'}$ and $M_j^{2-2\sigma'}\le1+(NX)^{2-2\sigma'}$.

\emph{Together.} With $|f|^2\le2|\sum_jP_j|^2+2|R|^2$ and $N^{2-2\sigma'}\le1+(NX)^{2-2\sigma'}$,
\[
\int_{T/2}^{3T}|f|^2\le4J^2L^5\bigl(3TX^{1-2\sigma'}+8+8(NX)^{2-2\sigma'}\bigr)+5(k+5)^2\bigl(1+(NX)^{2-2\sigma'}\bigr).
\]
For $T$ large, $J\le3\log T$ and $L\le4\log T$, which gives the claim.
\end{proof}

At $\sigma'=\frac12$ the bound is about $NX$, at $\sigma'=1$ about $\frac TX+1$: the longer the mollifier, the better it works to the right, and the worse to the left. Carlson's exponent comes from the best compromise.

\begin{theorem}[Carlson, 1920]\label{thm:carlson}
For every $\varepsilon>0$ there is $C(\varepsilon)$ such that for $\frac12\le\sigma\le1$ and $T\ge2$,
\[
N(\sigma,T)\le C(\varepsilon)\,T^{(4\sigma+\varepsilon)(1-\sigma)}(\log T)^{13}.
\]
\end{theorem}

\begin{proof}
We may take $\varepsilon\le1$. The main step is the bound
\[
N(\sigma,2T)-N(\sigma,T)\le C'(\varepsilon)T^{(4\sigma+\varepsilon)(1-\sigma)}(\log T)^{12}\qquad(T\ge T_2(\varepsilon)).\tag{$*$}
\]

\emph{Near the line.} Let $\eta=\frac1{\log T}$. If $\sigma<\frac12+\eta$, then $4\sigma(1-\sigma)=1-(2\sigma-1)^2\ge1-\frac4{\log^2T}$, so $T\le e^4T^{4\sigma(1-\sigma)}$, and ($*$) follows from $N(2T)-N(T)=O(T\log T)$. From now on $\frac12+\eta\le\sigma\le1$.

\emph{The set-up.} Let $\sigma_0=\sigma-\eta\ge\frac12$, $X=T^{2\sigma-1}\in[1,T]$, $f=\zeta M_X-1$ and $h=1-f^2$ as above. Let $c=3+\frac{\log T}{\log2}$, $H=2\sqrt{\eta(c-1)}$, and cover $\left]T,2T\right]$ by the windows $\left]t_k-\frac H2,t_k+\frac H2\right]$, $t_k=T+(k+\frac12)H$, $0\le k<K=\lceil\frac TH\rceil$. Let $n_k$ be the number of zeros of $\zeta$ with $\beta\ge\sigma$ in the $k$-th window, and $D_k$ the closed disc of centre $a_k=c+it_k$ and radius $R=c-\sigma_0$. For $T$ large, $c\le2\log T$, $2\le R\le c$, and since $H^2=\frac{4(c-1)}{\log T}$ lies between $\frac4{\log 2}$ and $8$, $2<H\le3$.

\emph{The centres.} At $\operatorname{Re}s=c>1$, $f(s)=\sum_{n\ge2}b_nn^{-s}$, so $|f(a_k)|\le\sum_{n\ge2}d(n)n^{-c}=\zeta(c)^2-1$. Comparing with an integral, $\zeta(c)-1\le2^{-c}+\int_2^\infty u^{-c}\dd u=2^{-c}(1+\frac2{c-1})\le2^{1-c}=\frac1{4T}$, and $\zeta(c)+1\le3$. So $|f(a_k)|\le\frac3{4T}$, $|h(a_k)|\ge1-|f(a_k)|^2$, and
$-\log|h(a_k)|\le2|f(a_k)|^2\le\frac2{T^2}$.

\emph{The geometry.} A zero $\rho$ in the $k$-th window with $\beta\ge\sigma$ satisfies $|\rho-a_k|\le r'$, where $r'^2=(c-\sigma)^2+\frac{H^2}4$. Since $\sigma,\sigma_0\le1$,
\[
R^2-r'^2=\eta(2c-\sigma-\sigma_0)-\frac{H^2}4\ge2\eta(c-1)-\eta(c-1)=\eta(c-1),
\]
so $R-r'=\frac{R^2-r'^2}{R+r'}\ge\frac{\eta(c-1)}{2c}\ge\frac\eta3$. Let $r''=\frac{R+r'}2\ge\frac R2\ge1$.

\emph{Jensen.} The disc $D_k$ lies in $\operatorname{Re}s\ge\sigma_0$ and $|\operatorname{Im}s-t_k|\le R$, far from the pole at $1$, so $h$ is holomorphic near it. For every $r\in[r'',R]$ except the finitely many radii of circles through zeros of $h$, Jensen's formula, in which all terms are $\ge0$ and the $n_k$ zeros above contribute at least $\log\frac r{r'}\ge\log\frac{r''}{r'}\ge1-\frac{r'}{r''}=\frac{R-r'}{2r''}\ge\frac{R-r'}{2R}$ each, gives
\[
n_k\frac{R-r'}{2R}\le\frac1{2\pi}\int_0^{2\pi}\log^+|h(a_k+re^{i\theta})|\dd\theta+\frac2{T^2}.
\]
Average over $r\in[r'',R]$, an interval of length $\frac{R-r'}2$. In polar coordinates $\dd A=r\dd r\dd\theta$ and $r\ge r''\ge1$, so $\int_{r''}^R\int_0^{2\pi}\log^+|h|\dd\theta\dd r\le I_k:=\iint_{D_k}\log^+|h|\dd A$. Hence, with $R-r'\ge\frac\eta3$,
\[
n_k\le\frac{2R}{R-r'}\Bigl(\frac{I_k}{\pi(R-r')}+\frac2{T^2}\Bigr)\le\frac{6R}\eta\Bigl(\frac{I_k}\eta+\frac2{T^2}\Bigr).
\]

\emph{Summing over the windows.} A point $w$ lies in $D_k$ only if $|\operatorname{Im}w-t_k|\le R$, which happens for at most $\frac{2R}H+1\le R+1$ values of $k$. All the discs lie in the rectangle $\sigma_0\le\operatorname{Re}w\le2c$, $\frac T2\le\operatorname{Im}w\le3T$. Since $\log^+|h|\le|f|^2$ and $K\le T$,
\[
N(\sigma,2T)-N(\sigma,T)\le\sum_kn_k\le\frac{6R}\eta\Bigl(\frac{R+1}\eta\int_{\sigma_0}^{2c}\int_{T/2}^{3T}|f(\sigma'+it)|^2\dd t\dd\sigma'+\frac2T\Bigr).
\]
With $R\le2\log T$, $R+1\le3\log T$, $2c\le4\log T$ and $\frac1\eta=\log T$, this is at most $144(\log T)^5\max_{\sigma_0\le\sigma'\le2c}\int_{T/2}^{3T}|f(\sigma'+it)|^2\dd t+1$.

\emph{The mean square.} By Proposition~\ref{prop:meansq}, the maximum is at most
\[
C(\varepsilon)(\log T)^7\bigl(TX^{1-2\sigma_0}+(NX)^{2-2\sigma_0}+1\bigr),
\]
and this bound decreases as $\sigma'$ increases. Now $X^{2\eta}\le T^{2/\log T}=e^2$, and $NX\le T^4$ for $T$ large, so $(NX)^{2\eta}\le e^8$. So we may replace $\sigma_0$ by $\sigma$, at the cost of a factor $e^8$. Finally
\[
TX^{1-2\sigma}=T^{1-(2\sigma-1)^2}=T^{4\sigma(1-\sigma)},\qquad NX\le25T^{2\sigma+\varepsilon/2},\quad(NX)^{2-2\sigma}\le25T^{(4\sigma+\varepsilon)(1-\sigma)},
\]
which proves ($*$).

\emph{Dyadic sum.} For $T\ge4T_2$, let $J\ge1$ be the integer with $T_2\le T2^{-J}<2T_2$. Apply ($*$) at $T2^{-j}$, $1\le j\le J$: each term is at most $C'(\varepsilon)T^{(4\sigma+\varepsilon)(1-\sigma)}(\log T)^{12}$, there are $J\le\log_2T$ of them, and $N(\sigma,T2^{-J})\le N(2T_2)$ is a constant. For $2\le T<4T_2$, $N(\sigma,T)\le N(4T_2)$ is a constant too.
\end{proof}

The exponent $4\sigma(1-\sigma)$ is exactly Carlson's, and it is a parabola: $1$ at $\sigma=\frac12$, where it says nothing beyond $N(\sigma,T)\le N(T)$, and $0$ at $\sigma=1$. Two consequences follow at once.

\begin{corollary}[Bohr and Landau, 1914]\label{cor:bohrlandau}
For every $\delta\in\left]0,\frac12\right[$, the number of zeros with $0<\gamma\le T$ and $|\beta-\frac12|\ge\delta$ is $O_\delta(T^{1-2\delta^2}(\log T)^{13})$. So almost all zeros, in proportion, lie in the strip $|\beta-\frac12|<\delta$.
\end{corollary}
\begin{proof}
By the symmetry $\rho\mapsto1-\overline\rho$, which keeps $\gamma$ and maps $\beta$ to $1-\beta$, the count is at most $2N(\frac12+\delta,T)$. Apply Theorem~\ref{thm:carlson} with $\sigma=\frac12+\delta$ and $\varepsilon=4\delta^2$: the exponent is $1-4\delta^2+4\delta^2(\frac12-\delta)\le1-2\delta^2$. Compare with $N(T)\sim\frac T{2\pi}\log T$.
\end{proof}

Bohr and Landau had the bound $O(T)$; Carlson's is a power smaller. Neither says anything about zeros \emph{on} the line: that is Hardy's theorem of 1914, in Chapter~14.

\begin{corollary}[primes in short intervals]\label{cor:hoheisel}
For every $\theta>\frac34$, $\psi(x+x^\theta)-\psi(x)\sim x^\theta$, and the number of primes in $\left]x,x+x^\theta\right]$ is asymptotic to $\frac{x^\theta}{\log x}$. In particular $p_{n+1}-p_n\le p_n^\theta$ for all large $n$.
\end{corollary}
\begin{proof}
Choose $\varepsilon>0$ with $1-\frac1{4+\varepsilon}<\theta$. Since $4\sigma\le4$, Theorem~\ref{thm:carlson} gives the hypothesis of Theorem~\ref{thm:ingham} with $A=4+\varepsilon$ and $B=13$, hence the first claim. Let $h=x^\theta$. The prime powers $p^j$, $j\ge2$, in $\left]x,x+h\right]$ are $m^j$ with
$2\le j\le\log_2(2x)$, and for each $j$ there are at most $(x+h)^{1/j}-x^{1/j}+1\le\frac{h}{2\sqrt x}+1$ of them (the mean value theorem, with $\frac1jx^{1/j-1}\le\frac1{2\sqrt x}$). They contribute at most $\log(2x)\log_2(2x)(\frac h{2\sqrt x}+1)=o(h)$ to $\psi(x+h)-\psi(x)$, since $\theta>\frac12$. So the primes in
the interval have $\sum\log p\sim h$, and each $\log p$ lies between $\log x$ and $\log 2x$. For the last claim take $x=p_n$: for $n$ large the interval $\left]p_n,p_n+p_n^\theta\right]$ contains a prime.
\end{proof}

\section{The race for the exponent}

\begin{figure}[ht]
\centering
\fig{ch01-density}
\caption{Bounds $N(\sigma,T)\ll T^{e(\sigma)+o(1)}$ for the number of zeros with real part at least $\sigma$: the exponent $e(\sigma)$ against $\sigma$. Carlson's parabola is the bound proved in this chapter. Ingham's curve (1940) uses the fourth moment of $\zeta$, and Huxley's straight line, valid for all $\sigma$ but useful only where it lies below
the others, comes from large values of Dirichlet polynomials. The density hypothesis, the dotted line, is still open. The dashed line is the trivial bound $N(\sigma,T)\le N(T)=T^{1+o(1)}$. Theorem~\ref{thm:ingham} turns the slope $A$ of a line $A(1-\sigma)$ lying above a curve into primes in all intervals of length $x^\theta$, $\theta>1-\frac1A$.}
\label{fig:density}
\end{figure}

Theorem~\ref{thm:ingham} reduces short intervals to a single number: the least $A$ such that $N(\sigma,T)\ll T^{A(1-\sigma)+o(1)}$ for all $\sigma$. Carlson's parabola lies below the line $4(1-\sigma)$, and touches it at $\sigma=1$, which is where the danger was. Ingham's 1937 paper used, instead of mean values alone, a bound on the critical line: if $\zeta(\frac12+it)=O(|t|^{c})$, then every $\theta>\frac{1+4c}{2+4c}$ works, and the value $c=\frac16+\varepsilon$, from the method of van der Corput, Hardy and Littlewood, gives $\theta>\frac58$. In 1940 he added the
bound $T^{\frac{3(1-\sigma)}{2-\sigma}}$, up to logarithms, from the fourth moment $\int_0^T|\zeta(\frac12+it)|^4\dd t=O(T\log^4T)$ that he had found in 1926; it is better than Carlson's for every $\sigma>\frac12$, though its slope $3$ at $\sigma=1$ only gives $\theta>\frac23$. Martin Huxley in 1972 proved $A=\frac{12}5$, hence $\theta>\frac7{12}$, by counting the large values of Dirichlet polynomials. The
\emph{density hypothesis}, $A=2$, would give every $\theta>\frac12$ (Exercise~\ref{ex:dh}): almost as good as the Riemann hypothesis for this question, which says that the zeros may stray from the line, provided they are rare enough. The present record for primes in short intervals, by Roger Baker, Glyn Harman and J\'anos Pintz in 2001, is $\theta=0.525$; it combines
zero density with sieve methods that only find a positive proportion of the expected primes, not an asymptotic. In 2024 Larry Guth and James Maynard proved $A=\frac{30}{13}$, the first improvement on Huxley's slope in fifty years, giving the asymptotic for $\theta>\frac{17}{30}$. All of these are quoted, not proved, in this book.

\begin{figure}[ht]
\centering
\fig{ch01-gaps}
\caption{The record gaps between consecutive primes below $10^9$, computed by a segmented sieve: each dot is a gap longer than all earlier ones, at the prime $p$ where it starts, from $p=3$ (gap $2$) to $p=436\,273\,009$ (gap $282$). Both axes are logarithmic. The bound under the Riemann
hypothesis, of the order of $\sqrt p\log p$, and the unconditional $p^{0.525}$ are far above; $(\log p)^2$, Cram\'er's conjecture, stays above all the records from $p=23$ on.}
\label{fig:gaps}
\end{figure}

What is true? Figure~\ref{fig:gaps} shows that the proved bounds are nowhere near the data. Cram\'er's conjecture $p_{n+1}-p_n=O((\log p_n)^2)$ came from his random model of Volume~5, and is still believed in a slightly modified form; the best proved bound, even under the Riemann hypothesis, is still about $\sqrt{p_n}\log p_n$. Every
one of the results above is a statement about all intervals. For most intervals much more is known: Atle Selberg proved in 1943, under the Riemann hypothesis, that almost all intervals $\left]x,x+h(x)\right]$ contain primes as soon as $h(x)/(\log x)^2\to\infty$.

\section{Exercises}

\begin{exercise}[$\star$]\label{ex:carlsonnum}
Use Theorem~\ref{thm:carlson} to show that, for every $\varepsilon>0$, at most $O_\varepsilon(T^{0.36+\varepsilon})$ zeros with $0<\gamma\le T$ have $\beta\ge0.9$. What does Ingham's bound $T^{\frac{3(1-\sigma)}{2-\sigma}}$ give?
\end{exercise}

\begin{exercise}[$\star$]\label{ex:mvtsharp}
Let $M\ge1$ and $0<T\le\frac1{4\log2}$. Show that
$\int_0^T\bigl|\sum_{M<n\le2M}n^{-it}\bigr|^2\dd t\ge\cos^2(\frac14)\,TM'^2$, where $M'$ is the number of integers in $\left]M,2M\right]$. Why does this show that the term $8M(1+\log M)$ in Lemma~\ref{lem:mvt} cannot be replaced by $o(M)$?
\end{exercise}

\begin{exercise}[$\star$]\label{ex:pk}
Show that $P_2(u)=\frac{\{u\}^2-\{u\}}2+\frac1{12}$ and $P_3(u)=\frac{\{u\}^3}6-\frac{\{u\}^2}4+\frac{\{u\}}{12}$. Deduce from Lemma~\ref{lem:em} with $k=2$ that $|\zeta(1+it)|\le2\log|t|+3$ for $|t|\ge3$.
\end{exercise}

\begin{exercise}[$\star\star$]\label{ex:dh}
Assume the density hypothesis in the form $N(\sigma,T)\ll T^{2(1-\sigma)}(\log T)^B$ for $\frac12\le\sigma\le1$. Show that $\psi(x+x^\theta)-\psi(x)\sim x^\theta$ for every $\theta>\frac12$. Why can the method not go below $\frac12$, even if all zeros were on the line?
\end{exercise}

\begin{exercise}[$\star\star$]\label{ex:discrete}
Let $M\ge1$, $F(t)=\sum_{M<n\le2M}a_nn^{-it}$, and let $t_1<\dots<t_R$ be points of $[0,T]$ with $t_{r+1}-t_r\ge1$. Show that
\[
\sum_{r=1}^R|F(t_r)|^2\le\bigl(T+1+8M(1+\log M)\bigr)\bigl(1+2\log(2M)\bigr)\sum_{M<n\le2M}|a_n|^2.
\]
\emph{(This discrete mean value theorem, a form of Gallagher's lemma, is the starting point of Huxley's method: zeros of $\zeta$ far to the right force large values of a Dirichlet polynomial at well-spaced points.)}
\end{exercise}

\begin{exercise}[$\star\star\star$]\label{ex:density1}
Let $\frac34<\sigma\le1$. Show that the mollifier of length $X=T$ gives, in the proof of Theorem~\ref{thm:carlson}, $N(\sigma,2T)-N(\sigma,T)\ll_\varepsilon T^{(4+\varepsilon)(1-\sigma)}(\log T)^{12}$, with no improvement on Carlson's exponent near $\sigma=1$. Then show that no choice of $X$ in $[1,T]$ can make the bound of Proposition~\ref{prop:meansq} at $\sigma'=\sigma$ smaller than $T^{4\sigma(1-\sigma)}$, so that a better $A$ needs a new idea.
\end{exercise}

\section{Hints}

\begin{hint}{ex:carlsonnum}
Compute $4\sigma(1-\sigma)$ and $\frac{3(1-\sigma)}{2-\sigma}$ at $\sigma=0.9$.
\end{hint}

\begin{hint}{ex:mvtsharp}
For $0\le t\le T$ and $M<n\le2M$, the phase $t\log\frac nM$ lies in $[0,\frac14]$. Look at the real part of $M^{it}\sum n^{-it}$.
\end{hint}

\begin{hint}{ex:pk}
Differentiate, check the mean, and use the uniqueness of a periodic antiderivative with mean zero. For $\zeta(1+it)$, take $N=\lceil(|s|+2)^2\rceil$.
\end{hint}

\begin{hint}{ex:dh}
Redo the proof of Theorem~\ref{thm:ingham} with $A=2$. For the second question, look at the zeros with $\beta=\frac12$ in the sum $S$.
\end{hint}

\begin{hint}{ex:discrete}
For a continuously differentiable $G$ on an interval $I$ of length $1$ containing $t_0$, $|G(t_0)|\le\int_I|G|+\int_I|G'|$. Apply it to $G=F^2$ on disjoint intervals of length $1$ and use Lemma~\ref{lem:mvt} for $F$ and $F'$.
\end{hint}

\begin{hint}{ex:density1}
In the bound $TX^{1-2\sigma}+(NX)^{2-2\sigma}$, one term decreases and the other increases with $X$. Compare their exponents.
\end{hint}

\section{Solutions}

\begin{solution}{ex:carlsonnum}
\idea Plug in: the theorem is a formula in $\sigma$.
\writeup At $\sigma=0.9$, $4\sigma(1-\sigma)=0.36$. Given $\varepsilon>0$, Theorem~\ref{thm:carlson} with $\varepsilon$ replaced by $\varepsilon/2$ gives $N(0.9,T)\le C\,T^{0.36+\varepsilon/20}(\log T)^{13}$, which is $O(T^{0.36+\varepsilon})$ because $(\log T)^{13}=O(T^{\varepsilon/2})$. Ingham's exponent is
$\frac{3\cdot0.1}{1.1}=\frac3{11}\approx0.273$, smaller. Either way, compared with $N(T)\approx\frac T{2\pi}\log T$, the zeros to the right of $0.9$ are a vanishing minority: if there are any, there are fewer than $T^{0.37}$ of them up to height $T$, for $T$ large.
\end{solution}

\begin{solution}{ex:mvtsharp}
\idea When $t$ is small, the terms $n^{-it}$ all point in nearly the same direction, so they add up without cancellation; the mean square is then $M'^2$, much more than the diagonal $M'$.
\writeup For $0\le t\le T$ and $M<n\le2M$, $0<t\log\frac nM\le T\log2\le\frac14$. So $\operatorname{Re}\bigl(M^{it}n^{-it}\bigr)=\cos(t\log\frac nM)\ge\cos\frac14$, and $\bigl|\sum_nn^{-it}\bigr|=\bigl|M^{it}\sum_nn^{-it}\bigr|\ge M'\cos\frac14$. Integrating the square over $[0,T]$ gives the claim.

Suppose the lemma held with $T+g(M)$ in place of $T+8M(1+\log M)$, where $g(M)=o(M)$. With $a_n=1$ and $T=\frac1{4\log2}$, it would give $\cos^2(\frac14)TM'^2\le(T+g(M))M'$, that is $M'\le\frac{T+g(M)}{T\cos^2(1/4)}=o(M)$, false since $M'\ge M-1$. So a term of size at least a constant times $M$ is needed; only the logarithm can be removed (Montgomery and Vaughan).
\end{solution}

\begin{solution}{ex:pk}
\idea A periodic function of mean zero is determined by its derivative, so it is enough to check derivative and mean.
\writeup On $]0,1[$, $Q_2(u)=\frac{u^2-u}2+\frac1{12}$ has $Q_2'=u-\frac12=P_1$, is continuous on $\R$ when extended with period $1$ (since $Q_2(0)=Q_2(1)$), and $\int_0^1Q_2=\frac16-\frac14+\frac1{12}=0$. Two continuous antiderivatives of $P_1$ differ by a constant, and the means fix it: $P_2=Q_2$. Likewise
$Q_3(u)=\frac{u^3}6-\frac{u^2}4+\frac u{12}$ has $Q_3'=\frac{u^2}2-\frac u2+\frac1{12}=P_2$, $Q_3(0)=Q_3(1)=0$, and $\int_0^1Q_3=\frac1{24}-\frac1{12}+\frac1{24}=0$, so $P_3=Q_3$.

For $s=1+it$, $|t|\ge3$, take $k=2$ and $N=\lceil(|s|+2)^2\rceil$. Lemma~\ref{lem:em} gives $|\zeta(s)|\le\sum_{n\le N}\frac1n+\frac1{|t|}+\frac3{2N}\le1+\log N+\frac13+\frac1{12}$. Here $|s|\le|t|+1$, so $N\le(|t|+3)^2+1\le(2|t|)^2+1$ for $|t|\ge3$, and $\log N\le2\log|t|+2\log2+\frac1{36}$. In total $|\zeta(1+it)|\le2\log|t|+1+1.39+0.03+0.42\le2\log|t|+3$.
\end{solution}

\begin{solution}{ex:dh}
\idea Only two places in the proof of Theorem~\ref{thm:ingham} use $A$: the choice of $T$ against $x$, and the requirement $\frac{T^A}x\to0$ with a power saving.
\writeup Take $A=2$, $h=x^\theta$, $T=x^{1-\theta}(\log x)^3$ with $\frac12<\theta\le1$. The proof of Theorem~\ref{thm:ingham} applies word for word: $1-\theta<\frac12$ makes $N(T)x^{-1/2}\to0$, and $\delta=1-2(1-\theta)=2\theta-1>0$. So $\psi(x+x^\theta)-\psi(x)\sim x^\theta$.

Now the limit of the method. The bound $|\psi(x+h)-\psi(x)-h|\le hS+O(\frac{x\log^2x}T)$ needs $T\gg\frac{x\log^2x}h$ to control the error. If all zeros had $\beta=\frac12$, then $S=\sum_{|\gamma|<T}x^{-1/2}\approx\frac T\pi\log T\cdot x^{-1/2}$, which is $\gg\frac{\sqrt x\log^3 x}h$; it tends to $0$ only if $h\gg\sqrt x\log^3x$. Bounding each zero's term by $hx^{\beta-1}$ ignores the cancellation between zeros, and with that loss even the Riemann hypothesis only gives $\theta>\frac12$ (it gives $h=\sqrt x\log^3x$, as in Volume~5).
\end{solution}

\begin{solution}{ex:discrete}
\idea A function that is large at a point stays large nearby unless its derivative is large; so a value is controlled by the integrals of $|G|$ and $|G'|$ near the point, and both integrals are mean values.
\writeup For $G$ continuously differentiable on an interval $I$ of length $1$ and $t_0,u\in I$, $G(t_0)=G(u)+\int_u^{t_0}G'$, so $|G(t_0)|\le|G(u)|+\int_I|G'|$; integrating over $u\in I$ gives $|G(t_0)|\le\int_I|G|+\int_I|G'|$. Take $G=F^2$, so $|G'|=2|F||F'|$, and $I_r=[t_r-\frac12,t_r+\frac12]$, which are disjoint except for endpoints and lie in $[-\frac12,T+\frac12]$. Summing, then Cauchy--Schwarz,
\[
\sum_r|F(t_r)|^2\le\int_{-1/2}^{T+1/2}|F|^2+2\Bigl(\int_{-1/2}^{T+1/2}|F|^2\Bigr)^{1/2}\Bigl(\int_{-1/2}^{T+1/2}|F'|^2\Bigr)^{1/2}.
\]
Here $F'(t)=\sum_n(-i\log n)a_nn^{-it}$ is a Dirichlet polynomial on the same block, with coefficients of modulus at most $|a_n|\log(2M)$. Lemma~\ref{lem:mvt} on an interval of length $T+1$ gives $\int|F|^2\le V\sum|a_n|^2$ and $\int|F'|^2\le V\log^2(2M)\sum|a_n|^2$ with $V=T+1+8M(1+\log M)$. So the sum is at most $V(1+2\log(2M))\sum|a_n|^2$.
\end{solution}

\begin{solution}{ex:density1}
\idea The proof of Theorem~\ref{thm:carlson} bounds the zero count by $(\log T)^{12}(TX^{1-2\sigma}+(NX)^{2-2\sigma})$ for any mollifier length $X\in[1,T]$, since $X$ entered only through Proposition~\ref{prop:meansq} and the inequalities $1\le X\le T$; so the question is about this function of $X$.
\writeup Write $X=T^a$, $0\le a\le1$, and $N=T^{1+\varepsilon/2+o(1)}$. Then
\[
\begin{gathered}TX^{1-2\sigma}=T^{1-a(2\sigma-1)},\qq(NX)^{2-2\sigma}=T^{(1+a+\varepsilon/2)(2-2\sigma)+o(1)}.\end{gathered}
\] With $a=1$: $T^{2-2\sigma}$ and $T^{(4+\varepsilon)(1-\sigma)+o(1)}$; the second dominates, and the proof gives $T^{(4+\varepsilon)(1-\sigma)}(\log T)^{12}$, the same slope $4$ at $\sigma=1$ as Carlson's choice
$a=2\sigma-1$, since $4\sigma\to4$.

For any $a$, the larger of the two exponents (with $\varepsilon\to0$) is $\max\bigl(1-a(2\sigma-1),(1+a)(2-2\sigma)\bigr)$. The first decreases in $a$, the second increases, and they are equal when $1-2a\sigma+a=2+2a-2\sigma-2a\sigma$, that is $a=2\sigma-1$, where both equal $4\sigma(1-\sigma)$. So the maximum is at least $4\sigma(1-\sigma)$ for every $a\in[0,1]$. With mean values of length $N\approx T$ and a mollifier, Carlson's exponent is the end of the road; Ingham's improvements need new input about $\zeta$ on the critical line, and Huxley's needs large values.
\end{solution}

\section*{Sources and further reading}

\begin{itemize}
\item \emph{Short intervals.} G.~Hoheisel, ``Primzahlprobleme in der Analysis'', \emph{Sitzungsber.\ Preuss.\ Akad.\ Wiss.\ Berlin} (1930), 580--588. A.~E.~Ingham, ``On the difference between consecutive primes'', \emph{Quart.\ J.\ Math.\ Oxford} 8 (1937), 255--266; ``On the estimation of $N(\sigma,T)$'', \emph{Quart.\ J.\ Math.\ Oxford} 11 (1940), 291--292; ``Mean-value theorems in the theory of the Riemann zeta-function'', \emph{Proc.\ London Math.\ Soc.}\ 27 (1926), 273--300. M.~N.~Huxley, ``On the difference between consecutive primes'',
\emph{Invent.\ Math.}\ 15 (1972), 164--170. R.~C.~Baker, G.~Harman and J.~Pintz, ``The difference between consecutive primes, II'', \emph{Proc.\ London Math.\ Soc.}\ 83 (2001), 532--562. L.~Guth and J.~Maynard, ``New large value estimates for Dirichlet polynomials'', preprint, 2024.
\item \emph{Zero density.} H.~Bohr and E.~Landau, ``Ein Satz \"uber Dirichletsche Reihen mit Anwendung auf die $\zeta$-Funktion und die $L$-Funktionen'', \emph{Rend.\ Circ.\ Mat.\ Palermo} 37 (1914), 269--272. F.~Carlson, ``\"Uber die Nullstellen der Dirichletschen Reihen und der Riemannschen $\zeta$-Funktion'', \emph{Arkiv f\"or Mat., Astron.\ och Fys.}\ 15 (1920), no.~20.
E.~C.~Titchmarsh, \emph{The Theory of the Riemann Zeta-Function}, 2nd ed.\ revised by D.~R.~Heath-Brown, Oxford, 1986, Chapter~9, for Carlson's and Ingham's theorems with Littlewood's lemma in place of Jensen's formula.
\item \emph{Mean values.} H.~L.~Montgomery and R.~C.~Vaughan, ``Hilbert's inequality'', \emph{J.\ London Math.\ Soc.}\ 8 (1974), 73--82. H.~L.~Montgomery, \emph{Topics in Multiplicative Number Theory}, Lecture Notes in Math.\ 227, Springer, 1971: mean values, large values and zero density in one place.
\item \emph{The zero-free region.} I.~M.~Vinogradov, ``A new estimate of the function $\zeta(1+it)$'', \emph{Izv.\ Akad.\ Nauk SSSR Ser.\ Mat.}\ 22 (1958), 161--164 (in Russian). N.~M.~Korobov, ``Estimates of trigonometric sums and their applications'', \emph{Uspehi Mat.\ Nauk} 13 (1958), no.~4, 185--192 (in Russian). K.~Ford, ``Vinogradov's integral and bounds for the Riemann zeta function'',
\emph{Proc.\ London Math.\ Soc.}\ 85 (2002), 565--633.
\item \emph{Gaps.} H.~Cram\'er, ``On the order of magnitude of the difference between consecutive prime numbers'', \emph{Acta Arith.}\ 2 (1936), 23--46. A.~Selberg, ``On the normal density of primes in small intervals, and the difference between consecutive primes'', \emph{Arch.\ Math.\ Naturvid.}\ 47 (1943), 87--105. The record gaps of Figure~\ref{fig:gaps} agree with the table of maximal gaps in the OEIS, sequence A002386.
\end{itemize}
